% ==============================================================================
% 10_criteria7_facilities.tex — Criterion 7: Facilities and Infrastructure
% Score: 2
% ==============================================================================

\criterion{7}{สิ่งอำนวยความสะดวกและโครงสร้างพื้นฐาน}{Facilities and Infrastructure}

% ------------------------------------------------------------------------------
\criterionsub{7.1 The programme to show that the physical resources are adequate to support the learning and research of students.}

หลักสูตรวิศวกรรมคอมพิวเตอร์และปัญญาประดิษฐ์จัดการเรียนการสอน ณ คณะเทคโนโลยีอุตสาหกรรม มหาวิทยาลัยราชภัฏอุตรดิตถ์ โดยใช้สิ่งอำนวยความสะดวกหลักในอาคาร ICIT (สำนักวิทยบริการและเทคโนโลยีสารสนเทศ) และอาคาร~10 (GE) ดังตาราง~\ref{tbl:classrooms}

\begin{table}[H]
  \centering
  \caption{ห้องเรียนและพื้นที่จัดการเรียนการสอน}
  \label{tbl:classrooms}
  \begingroup
  \footnotesize
  \setlength{\tabcolsep}{3pt}
  \renewcommand{\arraystretch}{1.2}
  \begin{tabularx}{\linewidth}{|>{\raggedright\arraybackslash}X|c|c|c|}
    \hline
    \textbf{อาคาร/ชั้น} & \textbf{ห้องเรียน/ห้องปฏิบัติการ} & \textbf{โสตทัศนูปกรณ์ (ชุด)} & \textbf{โต๊ะ/เก้าอี้} \\
    \hline
    ICIT ชั้น 1 & Co-Working Space + ห้องประชุม & -- & 10/30 \\
    \hline
    ICIT ชั้น 2 & ห้องปฏิบัติการ 7 ห้อง & 7 & 320/320 \\
    \hline
    ICIT ชั้น 3 & ห้องเรียน 5 ห้อง & 5 & 200 \\
    \hline
    ICIT ชั้น 4 & ห้องประชุม 2 ห้อง & -- & -- \\
    \hline
    อาคาร 10 (GE) ชั้น 2, 3 & Smart Classroom 13 ห้อง & 13 & 370 \\
    \hline
    URU Co-Working Space & 1 ห้อง & -- & 10/30 \\
    \hline
  \end{tabularx}
  \endgroup
\end{table}

ห้องเรียนและห้องปฏิบัติการทุกห้องติดตั้งโสตทัศนูปกรณ์ครบชุด ได้แก่ โปรเจคเตอร์หรือ Interactive Display ไมโครโฟน ระบบเสียง และระบบอินเทอร์เน็ต สามารถรองรับการสอนแบบ face-to-face และ hybrid ได้ ห้องเรียนสามารถจองผ่านระบบออนไลน์ที่ \url{https://academic.uru.ac.th/appl/admin/check_room.asp} (อ้างอิง~\hyperlink{ref:AUNQA-7-1-1}{AUNQA-7-1-1} ถึง~7-1-3)

สำหรับหลักสูตรระดับบัณฑิตศึกษา ขนาดชั้นเรียน~4 คนในภาค~2/2568 ทำให้สามารถจัดการเรียนการสอนแบบ seminar-style ที่มีปฏิสัมพันธ์สูงระหว่างอาจารย์กับนักศึกษา โดยใช้ห้องเรียน/ห้องประชุมขนาดเล็กที่มีความยืดหยุ่น

% ------------------------------------------------------------------------------
\criterionsub{7.2 The programme to show that laboratory equipment and facilities are adequate, up-to-date, readily available, and effectively deployed.}

หลักสูตรใช้ห้องปฏิบัติการคอมพิวเตอร์ในอาคาร ICIT ชั้น~2 จำนวน~7 ห้อง รวมเครื่องคอมพิวเตอร์ทั้งสิ้น~277 เครื่อง ดังตาราง~\ref{tbl:lab_computers}

\begin{table}[H]
  \centering
  \caption{จำนวนเครื่องคอมพิวเตอร์ในห้องปฏิบัติการ ICIT ชั้น 2}
  \label{tbl:lab_computers}
  \begingroup
  \footnotesize
  \renewcommand{\arraystretch}{1.2}
  \begin{tabular}{|l|c|}
    \hline
    \textbf{ห้อง} & \textbf{จำนวนเครื่อง} \\
    \hline
    C201 & 65 \\
    \hline
    C202 & 25 \\
    \hline
    C203 & 34 \\
    \hline
    C206 & 50 \\
    \hline
    C207 & 30 \\
    \hline
    C209 & 73 \\
    \hline
    \textbf{รวม} & \textbf{277} \\
    \hline
  \end{tabular}
  \endgroup
\end{table}

คอมพิวเตอร์ในห้องปฏิบัติการได้รับการปรับปรุงอย่างต่อเนื่อง โดยในปีการศึกษา~2568 ได้ดำเนินการเปลี่ยนคอมพิวเตอร์ใหม่ใน~2 ห้อง เพื่อให้รองรับซอฟต์แวร์ด้าน AI/ML ที่ต้องการทรัพยากรคำนวณสูง เช่น TensorFlow, PyTorch และ Jupyter Notebook ได้อย่างมีประสิทธิภาพ นักศึกษาสามารถเข้าใช้ห้องปฏิบัติการได้ตามตารางเวลาที่กำหนดและจองผ่านระบบออนไลน์ ซึ่งแสดงให้เห็นว่าทรัพยากรห้องปฏิบัติการถูกจัดสรรอย่างเป็นระบบ (Effectively deployed)

นอกจากห้องปฏิบัติการทางกายภาพ หลักสูตรยังสนับสนุนการใช้ Cloud Computing Resources ผ่านบริการ Google Colab, Kaggle Notebooks และ AWS Educate ซึ่งช่วยให้นักศึกษาสามารถฝึกโมเดล AI/ML ที่ต้องการ GPU ได้โดยไม่จำกัดเฉพาะห้องปฏิบัติการ

% ------------------------------------------------------------------------------
\criterionsub{7.3 The programme to show that there is a digital library with adequate and up-to-date resources and facilities that are in line with ICT advancement.}

มหาวิทยาลัยราชภัฏอุตรดิตถ์ให้บริการห้องสมุดดิจิทัลครบวงจรผ่านระบบ OPAC (\url{https://opac.uru.ac.th}) และฐานข้อมูลออนไลน์ที่สนับสนุนโดย สป.อว. ครอบคลุมสาขาวิศวกรรมคอมพิวเตอร์ วิทยาศาสตร์ข้อมูล และปัญญาประดิษฐ์ ดังนี้ (อ้างอิง~\hyperlink{ref:AUNQA-7-3-1}{AUNQA-7-3-1} ถึง~7-3-4)

ฐานข้อมูลหลักที่สนับสนุนงานวิจัยด้าน AI และวิศวกรรมคอมพิวเตอร์ประกอบด้วย Springer Nature Link สำหรับวารสารและหนังสือด้านวิทยาการคอมพิวเตอร์และวิศวกรรมศาสตร์ และ ScienceDirect (Elsevier) สำหรับวารสารและงานวิจัยด้าน Computer Science, AI และ Engineering ฐาน Engineering Source ให้บริการเฉพาะด้านวิศวกรรมศาสตร์ ในขณะที่ ACS (American Chemical Society) และ Academic Search Ultimate สนับสนุนงานวิจัยสหสาขาวิชา นอกจากนี้ยังมี ASTSU และ EBSCO EDS Plus Full Text ที่ให้บทความวิจัยเต็มรูปแบบ Emerald Management สำหรับมุมมองด้าน Business/Management ของเทคโนโลยี และ CU-eLibrary ที่รวบรวม e-book ทั้งภาษาไทยและภาษาอังกฤษ

ฐานข้อมูลเหล่านี้ครอบคลุมทั้ง IEEE, ACM Digital Library (ผ่าน ACM) และ Springer ซึ่งเป็นแหล่งตีพิมพ์หลักสำหรับงานวิจัยด้าน Computer Engineering และ AI ทำให้นักศึกษาสามารถเข้าถึงวารสารและการประชุมวิชาการระดับนานาชาติที่เกี่ยวข้องกับหัวข้อวิทยานิพนธ์ได้อย่างสะดวก

% ------------------------------------------------------------------------------
\criterionsub{7.4 The programme to show that the IT systems adequately support teaching, learning, research and service of academic staff, support staff and students.}

มหาวิทยาลัยราชภัฏอุตรดิตถ์ดำเนินการ Digital Transformation อย่างต่อเนื่อง โดยมีระบบสารสนเทศครอบคลุมทุกมิติของการจัดการการศึกษา ดังตาราง~\ref{tbl:it_systems}

\begin{table}[H]
  \centering
  \caption{ระบบ IT สนับสนุนการเรียนการสอน วิจัย และบริการ}
  \label{tbl:it_systems}
  \begingroup
  \footnotesize
  \setlength{\tabcolsep}{3pt}
  \renewcommand{\arraystretch}{1.2}
  \begin{tabularx}{\linewidth}{|>{\raggedright\arraybackslash}X|>{\raggedright\arraybackslash}X|}
    \hline
    \textbf{ระบบ} & \textbf{URL / ประโยชน์} \\
    \hline
    URU LMS (Life \& Learn) & \url{https://lms.uru.ac.th} --- จัดการเรียนการสอนออนไลน์ \\
    \hline
    ระบบภาระงาน & \url{https://workload.uru.ac.th} --- บันทึก/ติดตามภาระงานอาจารย์ \\
    \hline
    URU Lifelong Learning & \url{https://lifelong.uru.ac.th} --- คอร์สออนไลน์ LLL \\
    \hline
    ระบบอาจารย์ที่ปรึกษา & \url{https://advisor.uru.ac.th} --- ติดตามนักศึกษา \\
    \hline
    URU Exam Center & -- --- ทดสอบทักษะดิจิทัล \\
    \hline
    E-Meeting & \url{https://meeting.uru.ac.th} --- ประชุมออนไลน์ \\
    \hline
    ระบบตรวจ Plagiarism & \url{https://app.akarawisut.com/} --- ตรวจความซ้ำซ้อนงานวิจัย \\
    \hline
    Email มหาวิทยาลัย & email@live.uru.ac.th --- สื่อสารองค์กร \\
    \hline
    ระบบสารสนเทศประธานหลักสูตร & academic.uru.ac.th/teacher\_chief --- ดูผลการเรียนภาพรวม \\
    \hline
    ระบบบริการนักศึกษา & academic.uru.ac.th/std\_service1/ --- ผลการเรียน/ใบแจ้งหนี้ \\
    \hline
    ระบบจองห้องเรียน & academic.uru.ac.th/appl/admin/check\_room.asp --- จองห้อง online \\
    \hline
  \end{tabularx}
  \endgroup
\end{table}

ระบบ IT เหล่านี้ตอบสนองความต้องการของอาจารย์ สายสนับสนุน และนักศึกษาครบถ้วน (อ้างอิง~\hyperlink{ref:AUNQA-7-4-1}{AUNQA-7-4-1} ถึง~7-4-15) โดยเฉพาะระบบ URU LMS และระบบอาจารย์ที่ปรึกษาที่ใช้งานจริงในการบริหารจัดการหลักสูตรบัณฑิตศึกษา

% ------------------------------------------------------------------------------
\criterionsub{7.5 The programme to show that the computer and network infrastructure is adequate and accessible to academic staff and students throughout the institution.}

มหาวิทยาลัยราชภัฏอุตรดิตถ์ติดตั้งโครงสร้างพื้นฐานเครือข่ายผ่านระบบ \textbf{URU Smart Plus} ที่ครอบคลุมทุกพื้นที่ภายในมหาวิทยาลัย โดยมีจุด WiFi รวม \textbf{795 จุด} รองรับการเชื่อมต่อทั้ง LAN และ WiFi สำหรับอาจารย์ บุคลากร และนักศึกษาทั่วมหาวิทยาลัย

ความครอบคลุมของโครงสร้างเครือข่ายทำให้นักศึกษาสามารถเข้าถึง Cloud-based AI Tools (Google Colab, GitHub, Jupyter) ได้ทุกที่ในมหาวิทยาลัย ซึ่งเป็นปัจจัยสำคัญสำหรับหลักสูตรด้าน AI/ML ที่ต้องการการเชื่อมต่ออินเทอร์เน็ตความเร็วสูงอย่างต่อเนื่อง ข้อมูลโครงสร้างเครือข่ายอ้างอิงจากรายงานประจำปีของสำนักวิทยบริการและเทคโนโลยีสารสนเทศ (สวท.) มหาวิทยาลัยราชภัฏอุตรดิตถ์

% ------------------------------------------------------------------------------
\criterionsub{7.6 The programme to show that it complies with environment, health, safety and accessibility standards in the institution.}

มหาวิทยาลัยราชภัฏอุตรดิตถ์ดำเนินการตามมาตรฐานด้านสิ่งแวดล้อม สุขภาพ ความปลอดภัย และการเข้าถึงของผู้พิการ โดยอ้างอิงมาตรฐานระดับชาติ ดังนี้

มหาวิทยาลัยดำเนินการครอบคลุมหลายมิติ ด้านสิ่งแวดล้อม โครงการ Green Office / Green University ดำเนินการลดการใช้พลังงาน ส่งเสริมสิ่งแวดล้อม และเพิ่มพื้นที่สีเขียว สอดคล้องกับแนวทาง ISO~14001 ด้านการจัดการสิ่งแวดล้อม ด้านความปลอดภัยอาคาร อาคาร ICIT ผ่านการตรวจสอบตามกฎหมายควบคุมอาคาร (พ.ร.บ.\ ควบคุมอาคาร พ.ศ.~2522 และฉบับที่~2) มีระบบสัญญาณเตือนไฟไหม้ ถังดับเพลิง บันไดหนีไฟ และระบบไฟฉุกเฉิน เสริมด้วยกล้องวงจรปิดทั่วมหาวิทยาลัยและเวรยาม~24 ชั่วโมง ด้านสุขภาพ มีบริการรักษาพยาบาลเบื้องต้นและประกันอุบัติเหตุนักศึกษา~24 ชั่วโมง ด้านการเข้าถึงของผู้มีความต้องการพิเศษ มีคอมพิวเตอร์และเครื่องพิมพ์สำหรับผู้พิการทางสายตา ชั้น~2 อาคาร สวท. ลิฟต์และทางลาดสำหรับรถเข็นในอาคารหลัก ตามแนวทาง พ.ร.บ.\ ส่งเสริมและพัฒนาคุณภาพชีวิตคนพิการ พ.ศ.~2550 ส่วนมาตรฐานความปลอดภัยห้องปฏิบัติการวิจัยนั้น ห้องปฏิบัติการ CPE\&AI เป็นห้องปฏิบัติการด้านซอฟต์แวร์และ AI ที่ไม่มีสารเคมีหรืออุปกรณ์อันตราย และหลักสูตรอยู่ระหว่างยกระดับการกำกับดูแลตามแนวทาง ESPReL เพื่อรองรับการขยายขอบเขตวิจัยในอนาคต

% ------------------------------------------------------------------------------
\criterionsub{7.7 The programme to show that the physical, social, and psychological environment is conducive to education and research as well as student well-being.}

มหาวิทยาลัยราชภัฏอุตรดิตถ์จัดสภาพแวดล้อมที่เอื้อต่อการเรียนรู้ วิจัย และคุณภาพชีวิตนักศึกษาในหลายมิติ

ด้านสภาพแวดล้อมทางกายภาพ มหาวิทยาลัยจัดพื้นที่การเรียนรู้ที่หลากหลาย ได้แก่ Co-Learning Space และ Co-Working Space ในอาคาร ICIT ชั้น~1 สำหรับทำงานร่วมกัน หอพักนักศึกษาสำหรับผู้ที่ต้องการพักอาศัยในพื้นที่ สิ่งอำนวยความสะดวกด้านนันทนาการและสุขภาพซึ่งประกอบด้วยสนามกีฬา สระว่ายน้ำ และศูนย์ออกกำลังกาย รวมถึงโรงอาหาร ร้านสะดวกซื้อ ร้านกาแฟ และธนาคารภายในมหาวิทยาลัย นอกจากนี้ยังมีบริการรถรับ-ส่งนักศึกษาระหว่างมหาวิทยาลัยหลักกับวิทยาเขตลำรางทุ่งกะโล่ ด้านสภาพแวดล้อมทางจิตสังคม มหาวิทยาลัยจัดตั้งศูนย์รัก(ษ์)ใจเพื่อให้บริการสุขภาพจิตและให้คำปรึกษา มีห้องปฐมพยาบาล และระบบอาจารย์ที่ปรึกษาที่สนับสนุนการพัฒนาทางวิชาการและส่วนตัวของนักศึกษา

สำหรับนักศึกษาระดับบัณฑิตศึกษาโดยเฉพาะ พื้นที่ Co-Working Space ในอาคาร ICIT ชั้น~1 เป็นสภาพแวดล้อมที่เหมาะสมสำหรับการทำวิทยานิพนธ์/สารนิพนธ์และการประชุมกับอาจารย์ที่ปรึกษา

% ------------------------------------------------------------------------------
\criterionsub{7.8 The programme to show that the competences of the facilities and infrastructure support staff are defined, assessed and communicated.}

มหาวิทยาลัยราชภัฏอุตรดิตถ์กำหนดแนวทางปฏิบัติงานและมาตรฐานการให้บริการสำหรับเจ้าหน้าที่สายสนับสนุนด้านสิ่งอำนวยความสะดวก ครอบคลุมงานหลัก~5 ประเภท ได้แก่ งานอาคารสถานที่ งานรักษาความสะอาด งานรักษาความปลอดภัย งาน IT และงานบริการวิชาการ

เจ้าหน้าที่ได้รับการพัฒนาทักษะที่เกี่ยวข้อง เช่น การดูแลระบบโสตทัศนูปกรณ์ การแก้ไขปัญหา IT เฉพาะหน้า และการใช้เทคโนโลยีดิจิทัล มีการติดตามและประเมินผลการให้บริการจากข้อมูลความพึงพอใจและข้อเสนอแนะของผู้ใช้บริการ ซึ่งนำไปสู่การปรับปรุงอย่างต่อเนื่อง


% ------------------------------------------------------------------------------
\criterionsub{7.9 The programme to show that the quality of facilities and infrastructure is evaluated and benchmarked for improvement.}

หลักสูตรและมหาวิทยาลัยดำเนินการประเมินคุณภาพสิ่งอำนวยความสะดวกและโครงสร้างพื้นฐานอย่างสม่ำเสมอ ผลการประเมินความพึงพอใจสิ่งอำนวยความสะดวก ปีการศึกษา~2568 ดังตาราง~\ref{tbl:facility_satisfaction}

\begin{table}[H]
  \centering
  \caption{ผลการประเมินความพึงพอใจสิ่งอำนวยความสะดวก (5 คะแนน)}
  \label{tbl:facility_satisfaction}
  \begingroup
  \footnotesize
  \renewcommand{\arraystretch}{1.2}
  \begin{tabular}{|l|c|}
    \hline
    \textbf{หมวด} & \textbf{คะแนน} \\
    \hline
    ห้องเรียน & 4.32 \\
    \hline
    ห้องปฏิบัติการ & 4.32 \\
    \hline
    เทคโนโลยี IT & 4.29 \\
    \hline
    แหล่งเรียนรู้ (ห้องสมุด/ฐานข้อมูล) & 4.37 \\
    \hline
    สิ่งสนับสนุนทั่วไป & 4.38 \\
    \hline
    \textbf{ภาพรวม} & \textbf{4.34} \\
    \hline
  \end{tabular}
  \endgroup
\end{table}

ผลการประเมินอยู่ในระดับดีทุกด้าน (ภาพรวม~4.34/5.0) การดำเนินการปรับปรุงที่เกิดขึ้นจากผลการประเมินในปีการศึกษา~2568 ได้แก่ การเปลี่ยนคอมพิวเตอร์ใหม่ใน~2 ห้องปฏิบัติการ เพื่อเพิ่มประสิทธิภาพการรองรับงานด้าน AI/ML (อ้างอิง~\hyperlink{ref:AUNQA-7-9-1}{AUNQA-7-9-1} ถึง~7-9-2)

ในด้านการเทียบเคียงมาตรฐาน มหาวิทยาลัยเปรียบเทียบสิ่งอำนวยความสะดวกกับมหาวิทยาลัยราชภัฏในภาคเหนือและมหาวิทยาลัยที่เปิดสอนหลักสูตรด้าน AI/Computer Engineering เพื่อนำแนวทางที่ดีมาปรับใช้ในการพัฒนาอย่างต่อเนื่อง

\begin{strengthbox}
หลักสูตรมีสิ่งอำนวยความสะดวกครบวงจร: ห้องปฏิบัติการคอมพิวเตอร์~277 เครื่องในอาคาร ICIT, WiFi~795 จุดครอบคลุมทั่วมหาวิทยาลัย, ฐานข้อมูลออนไลน์นานาชาติ~7 ฐาน (Springer, ScienceDirect, Engineering Source ฯลฯ), ระบบ IT สนับสนุนการจัดการ~10 ระบบ และผลประเมินความพึงพอใจสิ่งอำนวยความสะดวกเฉลี่ย~4.34/5.0 มีการเปลี่ยนคอมพิวเตอร์ใหม่ใน~2 ห้องปฏิบัติการในปีการศึกษา~2568 ยิ่งกว่านั้น ผลสำรวจความพึงพอใจนักศึกษา CPE\&AI $n=4$ ปีการศึกษา~2568 ให้คะแนนด้านสิ่งสนับสนุนการเรียนรู้~5.00/5.00 (สูงสุด) ทุกด้าน ยืนยันคุณภาพจากผู้ใช้จริง
\end{strengthbox}

\begin{weaknessbox}
ยังไม่มีหลักฐาน ESPReL หรือมาตรฐานความปลอดภัยห้องปฏิบัติการวิจัยเฉพาะสำหรับ CPE\&AI (R7.6); กรอบสมรรถนะบุคลากรสายสนับสนุนด้านสิ่งอำนวยความสะดวกอยู่ระหว่างพัฒนา; Social Lab เฉพาะสำหรับนักศึกษาบัณฑิตศึกษาในระยะแรกยังไม่มีพื้นที่เฉพาะ
\end{weaknessbox}

\begin{selfassessmentbox}{2}{หลักสูตรมีสิ่งอำนวยความสะดวกและโครงสร้างพื้นฐานที่ครบถ้วนและมีคุณภาพสูง ครอบคลุมห้องปฏิบัติการ WiFi ฐานข้อมูลดิจิทัล และระบบ IT ผลประเมินความพึงพอใจอยู่ในระดับดี (4.34/5.0) และมีการปรับปรุงอย่างต่อเนื่อง แต่ยังต้องพัฒนากรอบสมรรถนะสายสนับสนุนด้านสิ่งอำนวยความสะดวกและมาตรฐานความปลอดภัยห้องปฏิบัติการวิจัยให้ครบถ้วนในรอบถัดไป}
\end{selfassessmentbox}

\begin{evidencelist}
\evidence{AUNQA-7-1-1}{รายการห้องเรียนและห้องปฏิบัติการ อาคาร ICIT และอาคาร 10}
\evidence{AUNQA-7-3-1}{ฐานข้อมูลออนไลน์: Springer, ScienceDirect, Engineering Source (สนับสนุนโดย สป.อว.)}
\evidence{AUNQA-7-4-1}{ระบบ URU LMS: \url{https://lms.uru.ac.th}}
\evidence{AUNQA-7-5-1}{ระบบเครือข่าย URU Smart Plus: WiFi 795 จุด}
\evidence{AUNQA-7-6-1}{โครงการ Green University + ระบบรักษาความปลอดภัย}
\evidence{AUNQA-7-7-1}{สิ่งอำนวยความสะดวกสนับสนุนคุณภาพชีวิต: Co-Working Space, ศูนย์รัก(ษ์)ใจ}
\evidence{AUNQA-7-9-1}{ผลการประเมินความพึงพอใจสิ่งอำนวยความสะดวก ปีการศึกษา 2568 (ภาพรวม 4.34/5.0)}
\evidence{AUNQA-C6-STSAT-2568}{แบบสำรวจความพึงพอใจนักศึกษาและผู้มีส่วนได้ส่วนเสีย CPE\&AI ปีการศึกษา~2568 ($n=4$): คะแนนด้านสิ่งสนับสนุนการเรียนรู้ (Lab/ฐานข้อมูล/WiFi/ห้องเรียน) ทุกรายการ~5.00/5.00 --- student\_satisfaction\_survey\_2568.csv \href{https://docs.google.com/forms/d/1xpfb6AfYkbOkidJNh6ci95ux4ZLYzjvpUuqARYjhr5k/viewanalytics}{\texttt{[Google Report]}}}
\end{evidencelist}
